\begin{abstract}
Optical coatings used in gravitational-wave detectors such as Advanced LIGO consist of multiple thin layers (${\sim}40$). Their thicknesses must be precisely controlled to meet stringent optical specifications. However, deposition errors can still occur due to imperfect deposition tools, which, in turn, can cause finished mirrors to deviate from their target specifications. In this paper, we investigate whether these deviations can be corrected during manufacturing by measuring the deposited layers and re-optimizing the thicknesses of the remaining layers. We call this procedure rescue. For this project, we used computer simulations of a 40-layer Advanced LIGO-style mirror coating. We simulated 32 different mirror coating designs that satisfy the required transmission targets at 1064~nm and 532~nm. Our simulations show that the success of rescue highly depends on the precision of measurements. For example, for independent, layer-to-layer deposition errors and perfect measurements, the rescue of the top 4 layers raised the largest error that still lets 95\% of mirrors pass from 1.14\% to 1.71\% for the typical (median) design. This method improved the tolerable error in 30 out of 32 designs, and re-optimizing 6 to 8 layers gave higher pass rates still. With 1\% measurement error, we still got improvement in 25 out of 32 designs. However, when measurement error reached 2\%, it reduced performance in 25 designs. Rescuing only the mirrors predicted to fail reduced this harm but did not remove it. Across all the cases we tested, the rescue helped only when the measurement error was smaller than the deposition error. Under fully shared deposition errors, the rescue produced only modest improvements: the median pass rate at 1\% error increased from 66\% to 76\%, but no design achieved a 95\% pass rate. Even with 12 re-optimized layers, only 2 of the 32 designs reached it, and a global optimizer did not help. The full thermal-noise calculation showed a negligible change: the median changed by $-0.02\%$.

These results suggest that in-process re-optimization can improve coating tolerance to independent deposition errors, provided that layers are measured more accurately than they are deposited. However, its effectiveness decreases substantially with measurement uncertainty, and it does not adequately compensate for systematic errors shared across all layers.
\end{abstract}
